\documentclass[11pt,a4paper]{article}

% ── Packages ──────────────────────────────────────────────────────────────
\usepackage[utf8]{inputenc}
\usepackage[T1]{fontenc}
\usepackage{amsmath,amssymb}
\usepackage{booktabs}
\usepackage{enumitem}
\usepackage{geometry}
\usepackage{hyperref}
\usepackage{graphicx}
\usepackage{xcolor}
\usepackage{fancyhdr}
\usepackage{titlesec}
\usepackage{longtable}
\usepackage{array}

\geometry{margin=2.5cm}
\hypersetup{colorlinks=true, linkcolor=blue!60!black, urlcolor=blue!60!black, citecolor=blue!60!black}

\pagestyle{fancy}
\fancyhf{}
\fancyhead[L]{\small FootyStats --- Algorithms \& Statistical Methods}
\fancyhead[R]{\small \thepage}
\renewcommand{\headrulewidth}{0.4pt}

\titleformat{\section}{\large\bfseries\color{blue!60!black}}{}{0em}{}[\vspace{-0.5em}\rule{\textwidth}{0.4pt}]
\titleformat{\subsection}{\normalsize\bfseries}{}{0em}{}

% ── Document ──────────────────────────────────────────────────────────────
\begin{document}

\begin{center}
  {\LARGE\bfseries FootyStats Analytics Platform}\\[6pt]
  {\large Algorithms \& Statistical Methods Reference}\\[10pt]
  {\normalsize Prepared for presentation and explanation of the unified
   season analytics + xG model system}\\[4pt]
  {\small\today}
\end{center}

\vspace{1em}

\tableofcontents
\newpage

% ======================================================================
\section{Overview}
% ======================================================================

The FootyStats system comprises two complementary machine-learning pipelines
that share a single player--season CSV as input:

\begin{enumerate}
  \item \textbf{Season Analytics} --- a regression-based pipeline that
        engineers per-game statistics, scores players, predicts season goals,
        and projects next-season performance.
  \item \textbf{Expected Goals (xG) Model} --- a classification-based pipeline
        that generates synthetic shot-level events and trains a probabilistic
        model estimating the likelihood that any given shot becomes a goal.
\end{enumerate}

The remainder of this document catalogues every algorithm, distribution,
formula, and evaluation metric used in each stage so you can explain exactly
\emph{what} was used and \emph{why}.


% ======================================================================
\section{Data Simulation \& Synthetic Shot Generation}
% ======================================================================

Because the raw data is at the player--season aggregation level (one row per
player per season), we must synthesise shot-level events before training the
xG model.  The following statistical methods are used.

\subsection{Estimating Number of Shots (Heuristic Rate Model)}

For each player we estimate total shots in the season as:
\[
  n_{\text{shots}} = \max\!\bigl(\,\text{round}(\text{appearances}
    \times r_{\text{pos}}),\; n_{\min}\bigr)
\]
where $r_{\text{pos}}$ is a position-dependent shots-per-appearance rate:

\begin{center}
\begin{tabular}{lc}
  \toprule
  Position   & $r_{\text{pos}}$ \\
  \midrule
  Forward    & 2.5  \\
  Midfielder & 1.0  \\
  Defender   & 0.3  \\
  Goalkeeper & 0.05 \\
  \bottomrule
\end{tabular}
\end{center}

and $n_{\min}=3$ ensures every player contributes training signal.

\subsection{Shot Outcome Simulation --- Binomial Distribution}

Given $n$ estimated shots and $g$ observed goals for a player, we model
each shot as an independent Bernoulli trial:
\[
  \text{is\_goal}_i \sim \mathrm{Bernoulli}(p), \qquad
  p = \min\!\Bigl(\frac{g}{n},\; 0.95\Bigr)
\]
Implemented via NumPy's \texttt{rng.binomial(1, p, size=n)}.  The cap at
0.95 prevents degenerate probabilities for very small $n$.

\subsection{Shot Location --- Gaussian (Normal) Distribution}

Pitch coordinates $(x, y)$ for each shot are sampled from position-dependent
Gaussian distributions.  The pitch is 105\,m $\times$ 68\,m with the
attacking goal at $(105, 34)$.

\[
  x \sim \mathcal{N}(\mu_x^{\text{pos}},\;\sigma_x^{\text{pos}}), \qquad
  y \sim \mathcal{N}(34,\;\sigma_y^{\text{pos}})
\]

\begin{center}
\begin{tabular}{lccc}
  \toprule
  Position   & $\mu_x$ & $\sigma_x$ & $\sigma_y$ \\
  \midrule
  Forward    & 90  & 5  & 10 \\
  Midfielder & 83  & 8  & 14 \\
  Defender   & 77  & 10 & 18 \\
  Goalkeeper & 75  & 12 & 20 \\
  \bottomrule
\end{tabular}
\end{center}

Values are clamped to the pitch boundaries $[0, 105] \times [0, 68]$.

\subsection{Contextual Feature Simulation}

\begin{itemize}[nosep]
  \item \textbf{Match minute} --- Discrete Uniform distribution,
        $U\{1, 90\}$.
  \item \textbf{Score difference} --- Categorical distribution over
        $\{-2, -1, 0, +1, +2\}$ with probabilities
        $(0.10, 0.20, 0.40, 0.20, 0.10)$.
  \item \textbf{Under pressure} --- Bernoulli$(0.35)$.
  \item \textbf{Body part} --- Categorical: foot ($p=0.82$),
        head ($p=0.18$).
  \item \textbf{Shot type} --- Categorical: open play ($p=0.80$),
        set piece ($p=0.20$).
\end{itemize}

\subsection{Stratified Train / Test Split}

The shot dataset is split 80\%/20\% using \textbf{stratified sampling} on
the binary \texttt{is\_goal} label, ensuring the goal/non-goal ratio is
preserved in both partitions.  Implemented via scikit-learn's
\texttt{train\_test\_split} with \texttt{stratify=y}.


% ======================================================================
\section{Feature Engineering}
% ======================================================================

\subsection{Shot Distance --- Euclidean Distance}

The straight-line distance from the shot location $(x, y)$ to the goal
centre $(x_g, y_g) = (105, 34)$:
\[
  d = \sqrt{(x_g - x)^2 + (y - y_g)^2}
\]

\subsection{Shot Angle --- Inverse Tangent (atan2)}

The bearing angle from the shot location to the goal centre:
\[
  \theta = \arctan\!\Bigl(\frac{|y - y_g|}{\max(x_g - x,\;10^{-6})}\Bigr)
\]
This returns a value in $[0, \frac{\pi}{2}]$ radians.  Smaller angles
correspond to more central, closer shots (higher scoring probability).
The $\varepsilon = 10^{-6}$ floor prevents division by zero.

\subsection{Player Performance Score --- Weighted Linear Combination}

In the season analytics pipeline, a composite performance score is computed:
\[
  S = 3.0 \cdot \frac{\text{goals}}{\text{apps}}
    + 1.5 \cdot \frac{\text{assists}}{\text{apps}}
    + 1.0 \cdot \frac{\text{apps}}{\max(\text{apps})}
    - 0.2 \cdot (\text{yellows} + 3 \times \text{reds})
\]

\subsection{Label Encoding}

Player positions (\textit{Forward}, \textit{Midfielder}, \textit{Defender},
\textit{Goalkeeper}) are mapped to integers via scikit-learn's
\texttt{LabelEncoder} for the season regression model.

\subsection{One-Hot Encoding}

For the xG classification model, categorical features (\texttt{body\_part},
\texttt{shot\_type}, \texttt{player\_position}) are converted to binary
dummy columns via scikit-learn's \texttt{OneHotEncoder}.

\subsection{Z-Score Standardisation}

Numeric features (shot distance, shot angle, match minute, score
difference, under pressure) are standardised to zero mean and unit
variance:
\[
  z = \frac{x - \bar{x}}{s}
\]
where $\bar{x}$ and $s$ are the training-set mean and standard deviation.
Implemented via scikit-learn's \texttt{StandardScaler}.


% ======================================================================
\section{Season Analytics --- Regression Models}
% ======================================================================

Three regression algorithms predict the total number of season goals per
player.

\subsection{Ridge Regression (L2-Regularised Linear Regression)}

Ridge regression minimises the sum of squared residuals plus an L2 penalty:
\[
  \hat{\boldsymbol{\beta}} = \arg\min_{\boldsymbol{\beta}}
  \left\{
    \sum_{i=1}^{n}(y_i - \mathbf{x}_i^\top \boldsymbol{\beta})^2
    + \alpha \|\boldsymbol{\beta}\|_2^2
  \right\}
\]
where $\alpha = 1.0$ controls regularisation strength.  The closed-form
solution is:
\[
  \hat{\boldsymbol{\beta}} = (\mathbf{X}^\top\mathbf{X}
    + \alpha\,\mathbf{I})^{-1}\mathbf{X}^\top\mathbf{y}
\]

\textbf{Why used:} Simple, interpretable baseline.  L2 regularisation
prevents overfitting on the small dataset and handles multicollinearity.

\subsection{Random Forest Regressor (Bagging Ensemble)}

An ensemble of $B = 200$ decision trees, each trained on a bootstrap
sample of the data and a random subset of features:
\[
  \hat{y} = \frac{1}{B}\sum_{b=1}^{B} T_b(\mathbf{x})
\]
where each tree $T_b$ is limited to \texttt{max\_depth}$\,= 4$.
Randomness is injected at two levels: \emph{bagging} (row sampling with
replacement) and \emph{feature bagging} (random column subsets at each
split).

\textbf{Why used:} Captures non-linear feature interactions that Ridge
cannot.  Limiting depth prevents overfitting.

\subsection{Gradient Boosting Regressor (Boosting Ensemble)}

Builds an additive model of $M = 100$ shallow trees sequentially, each
fitting the \emph{residuals} of the previous ensemble:
\[
  F_m(\mathbf{x}) = F_{m-1}(\mathbf{x})
    + \eta \cdot h_m(\mathbf{x})
\]
where $\eta$ is the learning rate and $h_m$ is a regression tree of
\texttt{max\_depth}$\,= 3$.  The algorithm minimises mean squared error
by gradient descent in function space.

\textbf{Why used:} Often the top performer for tabular data.  Sequential
correction of errors typically yields lower bias than bagging.

\subsection{Model Selection --- Leave-One-Out Cross-Validation (LOO-CV)}

For a dataset of $n$ players, LOO-CV trains $n$ models, each leaving one
player out for validation:
\[
  \mathrm{MAE}_{\text{LOO}} = \frac{1}{n}\sum_{i=1}^{n}
    |y_i - \hat{y}_{-i}|
\]
The model with the lowest LOO-CV MAE is selected.

\textbf{Why used:} With only $\sim$89 players, $k$-fold CV would yield
very small folds.  LOO-CV uses the maximum possible training data at each
iteration, giving the most unbiased error estimate for small datasets.


% ======================================================================
\section{xG Model --- Classification Algorithms}
% ======================================================================

Three classification algorithms estimate $P(\text{goal}\mid\text{shot
features})$.

\subsection{Logistic Regression (Baseline)}

Models the log-odds of a goal as a linear function of features:
\[
  \log\frac{P(y=1 \mid \mathbf{x})}{1 - P(y=1 \mid \mathbf{x})}
  = \boldsymbol{\beta}^\top \mathbf{x}
\]
Solving for the probability via the sigmoid function:
\[
  P(y=1 \mid \mathbf{x}) = \sigma(\boldsymbol{\beta}^\top \mathbf{x})
  = \frac{1}{1 + e^{-\boldsymbol{\beta}^\top \mathbf{x}}}
\]
Parameters are estimated by \textbf{Maximum Likelihood Estimation} (MLE),
optimised using the \textbf{L-BFGS} quasi-Newton method.  An L2 penalty
is applied with regularisation strength $C \in \{0.1, 1.0, 10.0\}$
(inverse of $\alpha$).

\textbf{Why used:} Directly outputs calibrated probabilities (the sigmoid
is itself a proper scoring function).  Serves as the interpretable
baseline.

\subsection{Random Forest Classifier (Bagging Ensemble)}

An ensemble of $B$ decision trees for classification.  The predicted
probability is the fraction of trees that vote ``goal'':
\[
  P(y=1 \mid \mathbf{x}) = \frac{1}{B}\sum_{b=1}^{B}
    \mathbb{1}[T_b(\mathbf{x}) = 1]
\]
Hyperparameter search space:
\begin{itemize}[nosep]
  \item \texttt{n\_estimators} $\in \{200, 400\}$
  \item \texttt{max\_depth} $\in \{\text{None}, 5, 10\}$
  \item \texttt{min\_samples\_leaf} $\in \{1, 5\}$
\end{itemize}
Class weights are set to \texttt{balanced\_subsample} to handle the
class imbalance ($\sim$14\% goal rate).

\textbf{Why used:} Robust to overfitting, naturally handles non-linear
boundaries and feature interactions.

\subsection{Gradient Boosting Classifier (Boosting Ensemble)}

Sequentially builds trees that correct previous errors, optimising the
\textbf{binomial deviance} (log loss) as the loss function:
\[
  \mathcal{L} = -\frac{1}{n}\sum_{i=1}^{n}\bigl[
    y_i \log p_i + (1 - y_i)\log(1 - p_i)
  \bigr]
\]
Hyperparameter search space:
\begin{itemize}[nosep]
  \item \texttt{n\_estimators} $\in \{100, 200\}$
  \item \texttt{learning\_rate} $\in \{0.05, 0.1\}$
  \item \texttt{max\_depth} $\in \{2, 3\}$
\end{itemize}

\textbf{Why used:} Typically achieves the best predictive accuracy for
tabular classification tasks.  The sequential error-correction is
particularly effective for probability estimation.

\subsection{Hyperparameter Tuning --- Grid Search with $k$-Fold CV}

All xG models are tuned via exhaustive \textbf{Grid Search} with
$k = 5$-fold cross-validation.  The scoring metric is
\textbf{negative log loss} (lower is better), ensuring we optimise
for probability quality rather than just classification accuracy:
\[
  \text{score} = -\mathcal{L}_{\text{log loss}}
\]
Implemented via scikit-learn's \texttt{GridSearchCV}.

\subsection{Probability Calibration --- Isotonic Regression}

After selecting the best model, we apply \textbf{isotonic regression
calibration} via \texttt{CalibratedClassifierCV} on held-out validation
data.  This is a non-parametric method that fits a monotonically
non-decreasing step function mapping raw model outputs to calibrated
probabilities:
\[
  \hat{p}_{\text{calibrated}} = f_{\text{isotonic}}(\hat{p}_{\text{raw}})
\]
where $f$ is a piecewise-constant, non-decreasing function fitted by
minimising:
\[
  \min_{f} \sum_{i=1}^{n} (y_i - f(\hat{p}_i))^2
  \quad\text{s.t.}\quad f \text{ is non-decreasing}
\]
solved via the Pool Adjacent Violators (PAV) algorithm.

\textbf{Why used:} Tree-based models often produce poorly calibrated
probabilities.  For xG, calibration is critical --- we need
$P(\text{goal}) = 0.15$ to truly mean 15\% of such shots are scored.


% ======================================================================
\section{Evaluation Metrics}
% ======================================================================

\subsection{Mean Absolute Error (MAE) --- Season Goal Model}
\[
  \mathrm{MAE} = \frac{1}{n}\sum_{i=1}^{n}|y_i - \hat{y}_i|
\]
Measures the average absolute difference between actual and predicted
season goals.  Expressed in the same units as the target (goals).

\subsection{Log Loss (Binary Cross-Entropy) --- xG Model}
\[
  \mathcal{L}_{\text{log}} = -\frac{1}{n}\sum_{i=1}^{n}
  \bigl[y_i\log(\hat{p}_i) + (1-y_i)\log(1-\hat{p}_i)\bigr]
\]
The primary metric for the xG model.  Heavily penalises confident wrong
predictions (e.g.\ predicting $\hat{p} = 0.95$ when $y = 0$).
Lower is better; a perfect model achieves 0.

\subsection{ROC-AUC (Area Under the Receiver Operating Characteristic Curve)}

The ROC curve plots True Positive Rate vs.\ False Positive Rate across
all decision thresholds.  AUC summarises discrimination ability:
\[
  \mathrm{AUC} = \int_0^1 \mathrm{TPR}(t)\;\mathrm{d}\,\mathrm{FPR}(t)
\]
\begin{itemize}[nosep]
  \item $\mathrm{AUC} = 0.5$: no discrimination (random classifier)
  \item $\mathrm{AUC} = 1.0$: perfect discrimination
\end{itemize}

\subsection{Brier Score}
\[
  \mathrm{BS} = \frac{1}{n}\sum_{i=1}^{n}(\hat{p}_i - y_i)^2
\]
The mean squared error of the predicted probabilities.
Ranges from 0 (perfect) to 1 (worst).

\subsection{Precision and Recall (at threshold $\tau = 0.5$)}
\[
  \text{Precision} = \frac{\mathrm{TP}}{\mathrm{TP} + \mathrm{FP}},
  \qquad
  \text{Recall} = \frac{\mathrm{TP}}{\mathrm{TP} + \mathrm{FN}}
\]
These are threshold-dependent metrics reported at the standard $\tau=0.5$
decision boundary for completeness.

\subsection{Confusion Matrix}

A $2 \times 2$ table of predicted vs.\ actual classes:

\begin{center}
\begin{tabular}{l|cc}
             & Predicted No Goal & Predicted Goal \\ \hline
  Actual No Goal & TN & FP \\
  Actual Goal    & FN & TP \\
\end{tabular}
\end{center}

\subsection{Calibration Curve (Reliability Diagram)}

Plots predicted probability (binned) against observed goal frequency in
each bin.  A perfectly calibrated model falls on the diagonal $y = x$.
We use 10 quantile-based bins.

\subsection{Feature Importance}

\begin{itemize}[nosep]
  \item \textbf{Tree-based models} --- Gini importance (mean decrease in
        impurity) accumulated across all trees.  Accessed via
        \texttt{feature\_importances\_}.
  \item \textbf{Linear models} --- Absolute coefficient magnitudes
        $|\beta_j|$ after standardisation.  Accessed via
        \texttt{coef\_}.
\end{itemize}


% ======================================================================
\section{Next-Season Projections --- Linear Extrapolation}
% ======================================================================

Player projections assume per-game rates remain constant:
\[
  \widehat{\text{goals}}_{\text{next}} = f_{\text{model}}(\text{proj\_features})
\]
where the trained regression model predicts from projected feature values:
\begin{align*}
  \text{proj\_apps}_i &= \text{round}\!\bigl(\text{involvement}_i \times N_{\text{games}}\bigr) \\
  \text{proj\_assists}_i &= \text{assists\_per\_game}_i \times \text{proj\_apps}_i \\
  \text{proj\_yellow}_i &= \frac{\text{yellows}_i}{\text{apps}_i} \times \text{proj\_apps}_i
\end{align*}
with $N_{\text{games}} = 12$.

Team-level projections are aggregated via summation, and a year-over-year
trend is computed as:
\[
  \text{trend}\% = \frac{\text{proj\_goals}_{\text{team}} - \text{current\_goals}_{\text{team}}}
    {\text{current\_goals}_{\text{team}}} \times 100
\]


% ======================================================================
\section{Summary Table --- Algorithm to Output Mapping}
% ======================================================================

\begin{longtable}{>{\raggedright}p{4.5cm} p{4.5cm} p{5.5cm}}
  \toprule
  \textbf{Output / Statistic} & \textbf{Algorithm} & \textbf{Key Formula / Method} \\
  \midrule
  \endhead

  Shot outcomes (is\_goal) &
    Binomial / Bernoulli trials &
    $\text{is\_goal} \sim \text{Bernoulli}(g/n)$ \\[4pt]

  Shot location $(x, y)$ &
    Gaussian (Normal) distribution &
    $x \sim \mathcal{N}(\mu, \sigma^2)$ per position \\[4pt]

  Shot distance &
    Euclidean distance &
    $d = \sqrt{(x_g - x)^2 + (y_g - y)^2}$ \\[4pt]

  Shot angle &
    Inverse tangent (atan2) &
    $\theta = \arctan(|y - y_g| / (x_g - x))$ \\[4pt]

  Performance score &
    Weighted linear combination &
    $S = 3g/a + 1.5a/a + 1.0i - 0.2d$ \\[4pt]

  Numeric scaling &
    Z-score standardisation &
    $z = (x - \bar{x}) / s$ \\[4pt]

  Categorical encoding &
    One-Hot Encoding &
    Binary dummy columns \\[4pt]

  Season goals prediction &
    Ridge / RF / GB Regression &
    LOO-CV, MAE selection \\[4pt]

  xG probability &
    LogReg / RF / GB Classification &
    5-fold GridSearchCV, log loss \\[4pt]

  Probability calibration &
    Isotonic regression (PAV) &
    Non-parametric monotone mapping \\[4pt]

  Model accuracy (goals) &
    LOO-CV MAE &
    $\frac{1}{n}\sum|y_i - \hat{y}_{-i}|$ \\[4pt]

  Model accuracy (xG) &
    Log loss, ROC-AUC, Brier score &
    See Section 5 \\[4pt]

  Feature importance &
    Gini importance / $|\beta|$ &
    Model-dependent \\[4pt]

  Next-season projections &
    Linear extrapolation + ML predict &
    Rate $\times$ projected appearances \\

  \bottomrule
\end{longtable}


% ======================================================================
\section{Software \& Libraries}
% ======================================================================

\begin{itemize}[nosep]
  \item \textbf{Python 3.13}
  \item \textbf{scikit-learn} --- Ridge, RandomForest, GradientBoosting,
        LogisticRegression, StandardScaler, OneHotEncoder, GridSearchCV,
        CalibratedClassifierCV, LOO-CV, all evaluation metrics
  \item \textbf{NumPy} --- random number generation (Gaussian, Binomial,
        Uniform, Categorical), array operations
  \item \textbf{pandas} --- data manipulation, groupby aggregations
  \item \textbf{matplotlib} --- diagnostic plots (ROC, PR, confusion
        matrix, calibration curve, feature importance)
  \item \textbf{joblib} --- model serialisation
\end{itemize}


\vfill
\begin{center}
\rule{0.6\textwidth}{0.4pt}\\[4pt]
{\small\itshape Generated for the FootyStats Analytics Platform.
All algorithms implemented via scikit-learn with reproducible
random seed (42).}
\end{center}

\end{document}
